\documentclass[11pt]{article}
\usepackage[margin=1.1in]{geometry}
\usepackage{amsmath,amssymb,amsthm}
\usepackage{booktabs}
\usepackage[colorlinks=true,linkcolor=blue,citecolor=blue,urlcolor=blue]{hyperref}

\newtheorem{theorem}{Theorem}[section]
\newtheorem{lemma}[theorem]{Lemma}
\newtheorem{corollary}[theorem]{Corollary}
\newtheorem{proposition}[theorem]{Proposition}
\theoremstyle{remark}
\newtheorem{remark}[theorem]{Remark}

\newcommand{\NA}{N_A}
\title{Exact $6$-cut rigidity and small-order superconnectivity\\
for the $6$-regular case of Dirac's $k=4$ problem}
% TODO(user): fill in author name / affiliation before submission
\author{Alper Ferudun\thanks{\texttt{alper@mercurycodelab.com}. Code, certificates
and machine-checked Lean files accompany this note as ancillary files; development
history at \texttt{github.com/AlperTheKing}.}}
\date{June 2026}

\begin{document}
\maketitle

\begin{abstract}
Dirac asked in 1970 whether for every $k\ge 4$ there is a $k$-vertex-critical graph
without critical edges; Jensen settled all $k\ge 5$, and only $k=4$ remains open.
Following Skottov\'a and Steiner, call a graph $G$ a \emph{$(4,1)$-graph} if
$\chi(G)=4$, $\chi(G-v)=3$ for every vertex $v$, and $\chi(G-e)=4$ for every edge
$e$; they proved $\delta(G)\ge 6$ and $\lambda(G)\ge 6$ for every $(4,1)$-graph and
asked whether a $6$-regular $(4,1)$-graph exists. We prove two results about this
$6$-regular case. \textbf{Theorem A} (computational, double-verified): there is no
$6$-regular $4$-vertex-critical graph on $n\le 12$ or $n=14$ vertices; on $n=13$
there is exactly one, and its $13$ critical edges form a Hamilton cycle.
Consequently any $6$-regular $(4,1)$-graph has at least $15$ vertices.
\textbf{Theorem B}: in a $6$-regular $(4,1)$-graph, every $6$-edge-cut is either the
edge star of a vertex or has both shores of size at least $15$; consequently every
$6$-regular $(4,1)$-graph on at most $29$ vertices is super-$6$-edge-connected.
The proof of Theorem~B combines an exact classification of the $3\times 3$ cut
matrices of $6$-edge-cuts in $(4,1)$-graphs (exactly $21$ matrices, five types up to
row/column permutations) with four local necessary conditions on a hypothetical
shore, the strongest being a \emph{boundary-shortfall lemma} derived from
vertex-deleted colourings; a finite, exhaustively enumerated candidate space then
collapses. The unique near-miss is $K_{3,3,3}$ minus a rainbow perfect matching of
size~$3$, which survives all counting filters and is excluded by the
boundary-shortfall lemma alone. Key lemmas are machine-checked in Lean~4/Mathlib,
and all enumerations were performed by two independent implementations.
\end{abstract}

\section{Introduction}

A graph $G$ is \emph{$k$-vertex-critical} if $\chi(G)=k$ and $\chi(G-v)=k-1$ for
every vertex $v$. An edge set $R\subseteq E(G)$ is \emph{critical} if
$\chi(G-R)<\chi(G)$. Dirac (see Erd\H{o}s \cite{Er89}) asked in 1970 whether for
every $k\ge 4$ there exists a $k$-vertex-critical graph with no critical edge.
Jensen \cite{Je02} resolved all cases $k\ge 5$; the case $k=4$ remains open, the
only prior partial result being that of Jensen and Siggers \cite{JS12}, who
constructed $4$-vertex-critical graphs in which fewer than an $\varepsilon$-fraction
of the edges are critical. The problem is \#944 in Bloom's database of Erd\H{o}s
problems \cite{Bl944}.

Following Skottov\'a and Steiner \cite{SkSt25} we call a graph $G$ a
\emph{$(4,1)$-graph} (for short, a \emph{target}) if
\[
\chi(G)=4,\qquad \chi(G-v)=3\ \ \forall v\in V(G),\qquad \chi(G-e)=4\ \
\forall e\in E(G).
\]
Skottov\'a and Steiner proved that every $(4,1)$-graph satisfies $\delta(G)\ge 6$
and $\lambda(G)\ge 6$ (their Proposition~5.1) and posed as an explicit subproblem
(their Problem~5.2) whether a \emph{$6$-regular} $(4,1)$-graph exists --- the
extremal, sparsest possible case. This note contributes two verified results on
that subproblem.

\begin{theorem}[A]\label{thm:A}
There is no $6$-regular $4$-vertex-critical graph on $n\le 12$ or on $n=14$
vertices. On $n=13$ vertices there is exactly one $6$-regular $4$-vertex-critical
graph $G_{13}$ (given in graph6 form in Section~\ref{sec:computations}); its
critical edges number $13$ and form a Hamilton cycle. Consequently every
$6$-regular $(4,1)$-graph has at least $15$ vertices.
\end{theorem}

\begin{theorem}[B]\label{thm:B}
In every $6$-regular $(4,1)$-graph, every $6$-edge-cut is either the edge star of a
single vertex or has both shores of size at least $15$. Consequently every
$6$-regular $(4,1)$-graph on at most $29$ vertices is super-$6$-edge-connected.
\end{theorem}

Theorem~B is obtained by an \emph{exact finite obstruction method}: the equality
case of the Skottov\'a--Steiner cut bound forces every $6$-edge-cut of a target to
carry one of exactly $21$ cut matrices (five types up to row and column
permutations, Theorem~\ref{thm:cutmatrix}); this, together with three further local
necessary conditions --- the strongest being the boundary-shortfall lemma
(Lemma~\ref{lem:shortfall}) --- confines a hypothetical small shore to a finite,
explicitly enumerable family which then collapses entirely. We consider the method
itself, and the structural near-miss it uncovers ($K_{3,3,3}$ minus a rainbow
$3$-matching, Section~\ref{sec:nine}), to be of independent interest.

We are careful to separate known from new. The bounds $\delta\ge 6$, $\lambda\ge 6$
and the cut-averaging argument are due to Skottov\'a and Steiner \cite{SkSt25};
the exclusion of comparable nonadjacent vertices in vertex-critical graphs is
folklore. The equality-case classification (Theorem~\ref{thm:cutmatrix}), the
boundary-shortfall lemma, the shore-exclusion theorems, and the computational
closure of $n=13,14$ appear to be new. None of this comes close to resolving
Problem~5.2: a target could still be a super-$6$-edge-connected $6$-regular graph
on $n\ge 15$ vertices, and our results do not constrain that core case.

All computations were performed by two independent implementations (C\texttt{++}
and Python) with exact agreement, with generator-side completeness guards, and the
central counting lemmas are machine-checked in Lean~4/Mathlib; see
Section~\ref{sec:computations} and the ancillary files.

\section{Local colouring lemmas}\label{sec:local}

Throughout, colourings are proper maps to $\{1,2,3\}$ (not counted up to colour
permutation). For $v\in V(G)$ and a $3$-colouring $\varphi$ of $G-v$ write
$a_i(v,\varphi)=|\{u\in N(v):\varphi(u)=i\}|$.

\begin{lemma}[singleton lemma]\label{lem:singleton}
Let $G$ be $4$-vertex-critical, $v\in V(G)$, and $\varphi$ a $3$-colouring of
$G-v$. Then $a_i(v,\varphi)\ge 1$ for every colour $i$; and if $a_i(v,\varphi)=1$,
say with unique colour-$i$ neighbour $u$, then the edge $vu$ is critical:
$\chi(G-vu)=3$.
\end{lemma}

\begin{proof}
If $a_i=0$, extend $\varphi$ by $\varphi(v)=i$; this $3$-colours $G$,
contradicting $\chi(G)=4$. If $a_i=1$ with witness $u$, keep $\varphi$ on
$G-v$ and set $\varphi(v)=i$ in $G-vu$: the only monochromatic edge would be
$vu$, which has been deleted.
\end{proof}

\begin{corollary}\label{cor:delta6}
In a target, $a_i(v,\varphi)\ge 2$ for all $v,\varphi,i$; hence $\delta(G)\ge 6$,
and in a $6$-regular target every neighbourhood splits exactly $2+2+2$ under every
vertex-deleted colouring. \textup{(}This is the local form of
\cite[Prop.~5.1]{SkSt25}.\textup{)}
\end{corollary}

\begin{corollary}\label{cor:quant}
Every $4$-vertex-critical graph $G$ has at least
$\tfrac12\sum_{v}\max(0,\,6-d(v))$ critical edges.
\end{corollary}

\begin{proof}
Fix for each $v$ one colouring $\varphi_v$ of $G-v$. The three positive integers
$a_i(v,\varphi_v)$ sum to $d(v)$, so at least $\max(0,6-d(v))$ of them equal $1$;
each singleton yields a critical edge at $v$ by Lemma~\ref{lem:singleton}, and
each critical edge is counted at most twice.
\end{proof}

\begin{lemma}[folklore]\label{lem:comparable}
A $k$-vertex-critical graph contains no two nonadjacent vertices $u,v$ with
$N(u)\subseteq N(v)$.
\end{lemma}

\begin{proof}
Take a $(k-1)$-colouring $\varphi$ of $G-u$ and set $\varphi(u)=\varphi(v)$:
every neighbour of $u$ is a neighbour of $v$, so no conflict arises and
$\chi(G)\le k-1$.
\end{proof}

We will also use the following observation, which gives a very effective
computational prefilter.

\begin{lemma}\label{lem:nbhd}
In a $4$-vertex-critical graph $G$ on $n\ge 8$ vertices, every induced
neighbourhood $G[N(v)]$ is bipartite.
\end{lemma}

\begin{proof}
An odd cycle $C\subseteq N(v)$ together with $v$ forms a $4$-chromatic wheel on at
most $7$ vertices; any vertex $u$ outside it gives $\chi(G-u)\ge 4$.
\end{proof}

\section{Exact rigidity of $6$-edge-cuts}\label{sec:cuts}

Let $G$ be a target and $V(G)=A\sqcup B$ a nontrivial partition with cut
$F=E(A,B)$. Both $G[A]$ and $G[B]$ are $3$-colourable (each is contained in some
$G-x$). For $3$-colourings $\alpha$ of $G[A]$ and $\beta$ of $G[B]$ define the
\emph{cut matrix} $M=(m_{ij})$ by
$m_{ij}=|\{ab\in F:\alpha(a)=i,\ \beta(b)=j\}|$, and for $\pi\in S_3$ the diagonal
sum $D_\pi(M)=\sum_i m_{i\pi(i)}$, the number of monochromatic cut edges after
permuting the colours of $B$ by $\pi$.

\begin{lemma}\label{lem:Dpi}
In a target: $D_\pi(M)=0$ is impossible; if $D_\pi(M)=1$ then the unique
monochromatic cut edge is critical. Hence $D_\pi(M)\ge 2$ for all $\pi$, and since
$\tfrac16\sum_{\pi}D_\pi(M)=|F|/3$, every cut satisfies $|F|\ge 6$
\textup{(\cite[Prop.~5.1]{SkSt25})}, and for $|F|=6$,
\[
D_\pi(M)=2\qquad\text{for all }\pi\in S_3\text{ and all pairs }(\alpha,\beta).
\]
\end{lemma}

\begin{proof}
$D_\pi=0$ gives a $3$-colouring of $G$; $D_\pi=1$ gives a $3$-colouring of $G-e$
for the unique monochromatic cut edge $e$. Each cut edge lies on exactly two of
the six permutation diagonals, giving the average $|F|/3$; for $|F|=6$ the average
equals the minimum, forcing equality everywhere.
\end{proof}

\begin{theorem}[$21$-matrix classification]\label{thm:cutmatrix}
Let $M$ be a $3\times 3$ matrix with nonnegative integer entries, total sum $6$,
and $D_\pi(M)=2$ for all $\pi\in S_3$. Then
\[
m_{ij}=\frac{R_i+C_j-2}{3}
\]
where $R_i,C_j$ are the row and column sums; all $R_i$ are congruent
$\bmod\ 3$, likewise all $C_j$; and $M$ is, up to row and column permutations, one
of the five types
\[
\begin{pmatrix}2&2&2\\0&0&0\\0&0&0\end{pmatrix},\
\begin{pmatrix}1&1&1\\1&1&1\\0&0&0\end{pmatrix},\
\begin{pmatrix}0&0&1\\0&0&1\\1&1&2\end{pmatrix},\
\begin{pmatrix}2&0&0\\2&0&0\\2&0&0\end{pmatrix},\
\begin{pmatrix}1&1&0\\1&1&0\\1&1&0\end{pmatrix},
\]
with orbit sizes $3,3,9,3,3$: exactly $21$ matrices in all. In particular the
row-sum vector is a permutation of $(6,0,0)$, $(3,3,0)$, $(4,1,1)$ or $(2,2,2)$.
\end{theorem}

\begin{proof}
Comparing two permutations that differ by a transposition and cancelling the
common entry gives $m_{rc}+m_{sd}=m_{rd}+m_{sc}$ for every $2\times 2$ submatrix;
hence $m_{pq}=m_{pj}+m_{iq}-m_{ij}$ for all $p,q,i,j$, and summing over all
$(p,q)$ yields $6=3R_i+3C_j-9m_{ij}$, i.e.\ the displayed formula. Integrality
forces $R_i+C_j\equiv 2 \pmod 3$ for all $i,j$, whence the congruence claims; the
nonnegative solutions of $\sum R_i=6$ with all $R_i$ congruent $\bmod\ 3$ are the
four listed vectors, and the case analysis over $(R,C)$ produces exactly the five
displayed types. The count $21$ is also machine-checked (Lean, by exhaustive
enumeration of the $3003$ weak compositions of $6$ into $9$ parts together with a
completeness lemma for the enumeration; see Section~\ref{sec:computations}).
\end{proof}

Since the row sums of $M(\alpha,\beta)$ do not depend on $\beta$, we obtain:

\begin{corollary}\label{cor:rowsum}
For every $3$-colouring $\alpha$ of a $6$-cut shore $G[A]$ in a target, the vector
$\bigl(\sum_{\alpha(v)=i} b(v)\bigr)_{i=1,2,3}$, where $b(v)$ is the number of cut
edges at $v$, is a permutation of $(6,0,0)$, $(3,3,0)$, $(4,1,1)$ or $(2,2,2)$.
\end{corollary}

\section{Shore constraints}\label{sec:shores}

From now on $G$ is a $6$-regular target and $A$ a shore of a nontrivial
$6$-edge-cut, $a=|A|\ge 2$, $H=G[A]$, $b(v)=6-d_H(v)$.

\begin{lemma}\label{lem:basic}
$e(H)=3a-3$; $H$ is connected; $0\le b(v)\le 5$ and $\sum_{v\in A}b(v)=6$.
\end{lemma}

\begin{proof}
Degree sum: $6a=2e(H)+6$. If $H$ had components $A_1,\dots,A_t$, the cuts
$\partial_G(A_i)$ are disjoint subsets of $\partial_G(A)$, each of size at least
$6$ by $\lambda(G)\ge 6$ \cite{SkSt25}, and their sizes sum to $6$; so $t=1$.
A vertex with $b(v)=6$ would be isolated in $H$.
\end{proof}

\begin{lemma}[boundary-shortfall lemma]\label{lem:shortfall}
For every $v\in A$ there exists a $3$-colouring $\psi$ of $H-v$ with
\[
\sum_{i=1}^{3}\max\bigl(0,\,2-c_i(\psi)\bigr)\ \le\ b(v),
\qquad\text{where } c_i(\psi)=|\{u\in \NA(v):\psi(u)=i\}|.
\]
\end{lemma}

\begin{proof}
Take a $3$-colouring $\varphi$ of $G-v$ and let $\psi=\varphi|_{A-v}$. By
Corollary~\ref{cor:delta6}, every colour appears at least twice on $N_G(v)$,
which is the disjoint union of $\NA(v)$ and the set $C_v$ of outside
cut-neighbours of $v$, $|C_v|=b(v)$. Writing $o_i$ for the number of vertices of
$C_v$ coloured $i$ we get $c_i+o_i\ge 2$, hence
$\max(0,2-c_i)\le o_i$, and summing over $i$ gives the bound since
$\sum_i o_i=b(v)$.
\end{proof}

For $v$ with $b(v)=0$ this says: \emph{$H-v$ must admit a colouring in which
$N(v)$ splits exactly $2+2+2$} --- full-degree vertices must be
``deletion-unfrozen''. Note also that for $u,v\in A$ nonadjacent with $b(u)=0$,
Lemma~\ref{lem:comparable} localizes: $\NA(u)\subseteq \NA(v)$ is impossible,
because then $N_G(u)=\NA(u)\subseteq N_G(v)$.

\section{Small shores}\label{sec:nine}

\begin{theorem}\label{thm:shore8}
No $6$-edge-cut in a target \textup{(}$6$-regular or not\textup{)} has a shore of
size $2\le a\le 8$.
\end{theorem}

\begin{proof}
By $\delta(G)\ge 6$, $e(H)\ge 3a-3$, while $H$ is $3$-colourable, so
$e(H)\le\lfloor a^2/3\rfloor$. For $2\le a\le 7$ these conflict
($\lfloor a^2/3\rfloor<3a-3$; this numeric fact is also Lean-checked). For $a=8$
both bounds equal $21$, forcing equality: $H$ is the Tur\'an graph
$K_{3,3,2}$ and every vertex of $A$ has $G$-degree exactly $6$. The two vertices
of the part of size $2$ then have all six edges inside $A$, so they are
nonadjacent vertices with equal neighbourhoods in $G$, contradicting
Lemma~\ref{lem:comparable}.
\end{proof}

\begin{theorem}\label{thm:shore913}
In a $6$-regular target, no $6$-edge-cut has a shore of size $9\le a\le 14$.
\end{theorem}

\begin{proof}[Proof (computer-assisted)]
By Lemma~\ref{lem:basic} a shore of size $a$ induces a connected graph on $a$
vertices with exactly $3a-3$ edges and maximum degree at most $6$. We enumerated
all such graphs with \texttt{geng} (nauty~2.8.9) and applied the necessary
conditions of Sections~\ref{sec:cuts}--\ref{sec:shores}: $3$-colourability;
Corollary~\ref{cor:rowsum} for \emph{every} $3$-colouring; the localized
Lemma~\ref{lem:comparable}; and Lemma~\ref{lem:shortfall} for every vertex. The
outcome:

\begin{center}
\begin{tabular}{rrrrrrr}
\toprule
$a$ & generated & not $3$-col. & fail Cor.~\ref{cor:rowsum} & fail
Lem.~\ref{lem:comparable} & fail Lem.~\ref{lem:shortfall} & survive\\
\midrule
9  & 729           & 711           & 9          & 8       & 1         & 0\\
10 & 18\,655        & 18\,345        & 197        & 86      & 27        & 0\\
11 & 696\,208       & 687\,377       & 6\,013      & 1\,300   & 1\,518     & 0\\
12 & 32\,833\,744    & 32\,484\,081    & 241\,863    & 27\,322  & 80\,478    & 0\\
13 & 1\,839\,349\,287 & 1\,822\,133\,664 & 11\,944\,366 & 788\,481 & 4\,482\,776 & 0\\
14 & 119\,236\,283\,370 & 118\,201\,012\,144 & 719\,933\,200 & 28\,146\,103 & 287\,191\,923 & 0\\
\bottomrule
\end{tabular}
\end{center}

No graph survives, proving the theorem. Verification methodology (independent
double implementation, generator completeness guards, per-graph certificates) is
described in Section~\ref{sec:computations}.
\end{proof}

\begin{remark}[the near-miss at $a=9$]
Exactly one graph survives all the \emph{counting} filters at $a=9$:
$H_9=K_{3,3,3}\setminus M_3$, where $M_3$ is a perfect rainbow matching of size
$3$ (one missing edge between each pair of colour classes, with six distinct
endpoints; graph6 \texttt{HEzftz\{}). $H_9$ has a unique $3$-colouring up to
colour permutation; its three full-degree vertices --- one in each class --- have
all neighbours inside $A$. For each such vertex $v$, all six colourings of
$H_9-v$ leave some colour with at most one occurrence on $N(v)$, so
Lemma~\ref{lem:shortfall} (with $b(v)=0$) excludes $H_9$. This is the only point
at $a\le 13$ where the boundary-shortfall lemma is essential for a graph that the
other filters miss, and we verified the kill additionally by brute force over all
$3^{8}$ assignments.
\end{remark}

\begin{proof}[Proof of Theorem~\ref{thm:B}]
A nontrivial $6$-edge-cut has shores of sizes $\ge 2$; Theorems~\ref{thm:shore8}
and \ref{thm:shore913} exclude sizes $2..14$, so both shores have size at least
$15$, forcing $n\ge 30$. For $n\le 29$ every $6$-edge-cut is therefore trivial,
i.e.\ a vertex star.
\end{proof}

\section{Computations, certificates, and machine-checked lemmas}
\label{sec:computations}

\paragraph{Theorem A.} All $6$-regular graphs on $n\le 14$ vertices were
enumerated with \texttt{geng -d6 -D6} (nauty 2.8.9) and classified by a C\texttt{++}
backtracking $3$-colourability checker into: $3$-colourable; $\chi\ge 4$ but not
vertex-critical; $4$-vertex-critical with a critical edge; target. Counts:

\begin{center}
\begin{tabular}{rrrrrr}
\toprule
$n$ & total & $3$-colourable & not vertex-critical & $4$-VC with critical edge &
target\\
\midrule
11 & 266 & 3 & 263 & 0 & 0\\
12 & 7\,849 & 50 & 7\,799 & 0 & 0\\
13 & 367\,860 & 849 & 367\,010 & 1 & 0\\
14 & 21\,609\,301 & 42\,667 & 21\,566\,634 & 0 & 0\\
\bottomrule
\end{tabular}
\end{center}

The $n=14$ run was executed twice with different residue-class chunkings
(mod~$110$ and mod~$73$) with identical aggregates, and the total matches the
known count of $6$-regular graphs on $14$ vertices. The unique $n=13$ graph is
\begin{center}
\texttt{L?bFFbw\textasciitilde B\{FwFw}
\end{center}
Its $13$ critical edges form a Hamilton cycle. Across its $13$ vertex-deleted
subgraphs it admits exactly $78$ proper $3$-colourings (as maps to $\{1,2,3\}$,
not modulo $S_3$), and \emph{none} of them splits the open neighbourhood
$2+2+2$: every colouring leaves a singleton class in $N(v)$, so
Lemma~\ref{lem:singleton} predicts critical edges at every vertex; the $156$
singleton predictions hit exactly the $13$ critical edges (no false positives).
Thus the unique smallest $6$-regular $4$-vertex-critical graph fails the target
conditions everywhere locally, not marginally.

\paragraph{Theorem B.} Candidate shores were generated by
\texttt{geng -c -D6 $a$ $E$:$E$} with $E=3a-3$; large cases ran as residue
classes with a generator-side completeness guard (each class must end with
\texttt{geng}'s \texttt{>Z} terminator line, and the generated count must equal
the classifier's count). The filter battery was implemented twice independently
(C\texttt{++} and Python) with identical counts at $a=9,10,11$; the $a=9$ kill was
additionally brute-forced. Per-graph certificates (graph6 plus the failed filter,
and for shortfall kills the killing vertex) are included for $a\le 11$ in the
ancillary files.

\paragraph{Machine-checked lemmas.} The following are verified in Lean~4 /
Mathlib (no \texttt{sorry}, no \texttt{native\_decide}; axioms at most
\texttt{propext}, \texttt{Classical.choice}, \texttt{Quot.sound}):
the recolouring core of Lemma~\ref{lem:singleton} (axiom-free); the count of
$21$ in Theorem~\ref{thm:cutmatrix}, via an enumeration of the $3003$ weak
compositions of $6$ into $9$ parts together with a proved completeness lemma for
the enumeration, and a bridge lemma stating the matrix form with
$\forall\pi\in S_3$ diagonal bounds; and the numeric inequality
$\lfloor a^2/3\rfloor<3a-3$ for $2\le a\le 7$ used in Theorem~\ref{thm:shore8}.

\section{Concluding remarks}

Theorem~B turns the Skottov\'a--Steiner edge-connectivity threshold into an exact
small-cut impossibility statement, but it says nothing about
super-$6$-edge-connected candidates, which is where Problem~5.2 now lives: a
$6$-regular target must have $n\ge 15$ (Theorem~A), and if $n\le 29$ it has no
nontrivial $6$-edge-cut at all (Theorem~B). The shore-exclusion machine is
open-ended --- each further size $a$ is a finite computation; the $a=14$ case
($119{,}236{,}283{,}370$ candidates, every residue class closed with a
generator-side completeness guard) completed in under a day on a workstation --- but we
do not currently see an argument that it terminates for all $a$: the row-sum
filter can be neutralized by small colour-forcing gadgets, and the burden falls
entirely on the boundary-shortfall lemma. Proving that
Corollary~\ref{cor:rowsum} forces enough colour rigidity for
Lemma~\ref{lem:shortfall} to fail at some full-degree vertex, for every $a$,
would upgrade Theorem~B to: \emph{every $6$-regular $(4,1)$-graph is
super-$6$-edge-connected}. We leave this as the natural next question.

\begin{thebibliography}{9}

\bibitem{Bl944} T.~Bloom, \emph{Erd\H{o}s problem \#944}, Online database of
Erd\H{o}s problems, \url{https://www.erdosproblems.com/944}.

\bibitem{Er89} P.~Erd\H{o}s, \emph{On some aspects of my work with Gabriel
Dirac}, in: Graph Theory in Memory of G.~A.~Dirac (Sandbjerg, 1985), Ann.
Discrete Math.\ \textbf{41}, North-Holland, 1989, 111--116.

\bibitem{Je02} T.~R.~Jensen, \emph{Dense critical and vertex-critical graphs},
Discrete Math.\ \textbf{258} (2002), 63--84.

\bibitem{JS12} T.~R.~Jensen and M.~Siggers, \emph{On a question of Dirac on
critical and vertex-critical graphs}, Sib.\ \`Elektron.\ Mat.\ Izv.\ \textbf{9}
(2012), 156--160.

\bibitem{SkSt25} E.~Skottov\'a and R.~Steiner, \emph{Critical edge sets in
vertex-critical graphs}, arXiv:2508.08703 (2025).

\end{thebibliography}

\end{document}
